\subsection{映射的喷射}

12.1.1. 设 \(x \in X\)，且 \(f, g\) 是在 \(x\) 的邻域内定义、取值于 \(Y\) 的两个连续映射。我们称 \(f\) 与 \(g\) 有阶数 \(\geqslant k\) 的接触于 \(x\) 处，是指 \(f(x) = g(x)\)，并且存在图表 \((U, \varphi, E)\)，它是 \(X\) 在 \(x\) 处的图表，以及图表 \((V, \psi, F)\)，它是 \(Y\) 在 \(f(x)\) 处的图表，使得映射 \(\psi \circ f \circ \varphi^{-1}\) 和 \(\psi \circ g \circ \varphi^{-1}\) 在 \(\varphi(x)\) 处有阶数 \(\geqslant k\) 的接触，并且在 E 中 \(\varphi(x)\) 的某邻域内定义、取值于 F。（1.1.2）。于是，这一性质对 \(X\) 在 \(x\) 处的所有图表以及 \(Y\) 在 \(f(x)\) 处的所有图表都成立。

12.1.2. 设 \(x \in \mathbf{X}\)，\(y \in \mathbf{Y}\)。在类为 \(\mathbf{C}^r\)、定义于 \(x\) 的某邻域、取值于 \(\mathbf{Y}\) 并将 \(x\) 映到 \(y\) 的映射集合中，关系“\(f\) 与 \(g\) 有阶数 \(\geqslant k\) 的接触”是等价关系；在此关系下 \(f\) 的等价类记作 \(j_x^k(f)\)，称为 \(k\) 阶的 \(f\) 的喷射，其源为 \(x\)，靶为 \(y\)。喷射 \(j\) 的源记作 \(s(j)\)，靶记作 \(b(j)\)。

从 X 到 Y 的 k 阶喷射集合（相应地，源为 x 的喷射集合、目标为 y 的喷射集合）记作 \(J^{k}(X, Y)\)（相应地记作 \(J_{x}^{k}(X, Y)\)、\(J^{k}(X, Y)_{y}\)），并置

\[
\mathrm{J} _ {x} ^ {k} (\mathrm{X}, \mathrm{Y}) _ {y} = \mathrm{J} _ {x} ^ {k} (\mathrm{X}, \mathrm{Y}) \cap \mathrm{J} ^ {k} (\mathrm{X}, \mathrm{Y}) _ {y}.
\]

12.1.3. 若 U 和 V 分别是 X 和 Y 的开子集，则按显然方式将 \(\mathrm{J}^{k}(\mathrm{U},\mathrm{V})\) 与 \(U \times V\) 在该映射下的逆像等同

\[
(s, b): \mathrm{J} ^ {k} (\mathrm{X}, \mathrm{Y}) \rightarrow \mathrm{X} \times \mathrm{Y}.
\]

12.1.4. 设 Z 是类 \(C^{r}\) 的流形，且 \((x, y, z) \in X \times Y \times Z\)。设 \(j \in \mathrm{J}_{x}^{k}(X, Y)_{y}\)，\(j' \in \mathrm{J}_{y}^{k}(Y, Z)_{z}\)。映射 \(f' \circ f\)，其中 \(f \in j\) 且 \(f' \in j'\)，在 x 处有相同的 k 阶喷射；此喷射称为 \(j'\) 与 j 的复合，记为 \(j' \circ j\)；有 \(s(j' \circ j) = s(j)\) 且 \(b(j' \circ j) = b(j')\)。若 T 是类 \(C^{r}\) 的流形且 \(j'' \in J_{z}^{k}(Z, T)\)，则有

\[
j ^ {\prime \prime} \circ (j ^ {\prime} \circ j) = (j ^ {\prime \prime} \circ j ^ {\prime}) \circ j.
\]

12.1.5. 设 \(j \in \mathrm{J}_x^k(\mathrm{X}, \mathrm{Y})\)，其中 \(x \in \mathrm{X}\)，并令 \(k'\) 为满足 \(0 \leqslant k' \leqslant k\) 的整数。对于喷射 \(j_x^{k'}(f)\)（\(f \in j\)），它们都相等；由此定义的喷射记作 \(r^{k, k'}(j)\)；映射 \(r^{k, k'}: \mathrm{J}^k(\mathrm{X}, \mathrm{Y}) \to \mathrm{J}^{k'}(\mathrm{X}, \mathrm{Y})\) 是满射。

12.1.6. 假设 \(r\) 等于 \(\infty\) 或 \(\omega\)。令 \(f\) 和 \(g\) 是类 \(C^{r}\) 的两个映射，在点 \(x \in X\) 的某邻域内定义并取值于 Y。若 \(j_x^m(f) = j_x^m(g)\) 对每个整数 \(m \geqslant 0\) 都成立，则称 \(f\) 与 \(g\) 在 \(x\) 处有无穷阶接触。由此得到一个等价关系，其等价类称为 \(x\) 处的无穷阶喷射；\(f\) 的无穷阶喷射记作 \(j_x^\infty(f)\) 或 \(j_x^\omega(f)\)。上述定义和结果不加修改地推广到无穷阶喷射。

